\section{The collective--transverse mechanism}\label{sec:network-mechanism}

This section isolates the mechanism behind
\cref{thm:transverse,cor:canonical-transfer,thm:one-gap}.  The point is not
that a synchronous orbit can be copied into \(N\) nodes; that observation
alone gives no uniform inverse and no control of nonsynchronous histories.
The useful fact is that the collective RFDE and every transverse network
direction admit estimates in norms whose constants do not grow with \(N\).

\subsection{Diameter contraction and the delay margin}

Write a variational solution along a complete synchronous trajectory as
\[
 (x,y)=(X\mathbf1,Y\mathbf1)+(x_\perp,y_\perp),
 \qquad \pi^Tx_\perp=\pi^Ty_\perp=0.
\]
For a complex vector \(z\), put
\(\operatorname{diam} z=\max_{i,k}|z_i-z_k|\).  This is a norm on
\(\Ker\pi^T\).  The Dobrushin identities imply
\begin{equation}
 \operatorname{diam}(Qz)\le\tau(Q)\operatorname{diam} z,\qquad
 \operatorname{diam}(B_jz)\le\frac12\operatorname{diam} z.
                                                               \label{eq:dobrushin-diam}
\end{equation}
No eigenbasis, reversibility, or simultaneous diagonalization is used.

Let
\[
 M(t)=\max\{\operatorname{diam} x_\perp(t),
             3\operatorname{diam} y_\perp(t)\}.
\]
On the strip \(|V(t)-1|\le5/2\), the current and delayed coefficients in
the leaky model give
\begin{align}
 \alpha_v&\ge 3(1-\tau(Q))-(1-\eps\kappa_1)-\frac13
          >0.16746,\\
 \alpha_w&\ge2(1-\tau(Q))+\eps-3\eps\ge0.6,\\
 \beta&\le\eps\{\kappa_1+3\kappa_3(5/2)^2\}=0.01955.
                                                               \label{eq:halanay-data}
\end{align}
Consequently
\begin{equation}
 D^+M(t)\le-\alpha M(t)
       +\beta\sup_{t-r\le s\le t}M(s),\qquad
 \alpha=\min\{\alpha_v,\alpha_w\}.                     \label{eq:halanay}
\end{equation}
The directed residual
\begin{equation}
 \alpha-\frac1{10}-\beta e^{r/10}
 >0.0076664505356400                                    \label{eq:delay-margin}
\end{equation}
is strict.  A first-crossing argument for
\(e^{t/10}M(t)\) proves
\[
 \norm{U_{\perp,N}(t,s)}_{1/10}\le e^{-(t-s)/10}.
\]
This proves \cref{eq:green-bound}'s homogeneous part in a norm adapted to
node oscillation.  It deliberately avoids a dimension-dependent
equivalence with a Euclidean product norm.

\subsection{A causal complete-line inverse}

The inverse in \cref{thm:transverse} is obtained without evolving an RFDE
backwards.  Given a bounded transverse forcing \(f\), let \(z^S\) be the
forward solution with zero retained history at time \(S\).  The forced
version of the same first-crossing argument gives
\begin{equation}
 \norm{z_t^S}_{1/10}
 \le\int_S^t e^{-(t-u)/10}\norm{f(u)}_\perp\,\dd u
 \le10\norm{f}_{\perp,\infty}.                         \label{eq:forced-pullback}
\end{equation}
If \(S_1<S_2<t\), the two solutions differ homogeneously after \(S_2\);
hence
\[
 \norm{z_t^{S_1}-z_t^{S_2}}_{1/10}
 \le10e^{-(t-S_2)/10}\norm f_{\perp,\infty}.
\]
The pullback limit \(S\to-\infty\) is therefore a unique bounded complete
solution and defines \(G_{\perp,N}\).  This proves
\(\norm{G_{\perp,N}}\le10\).  The construction is causal and remains valid
for nonperiodic complete synchronous trajectories as long as the voltage
strip holds.

\subsection{Exact transfer and structural transfer}

For the canonical synchronized Lin realization, all endpoint, phase, and
gap rows are collective.  Row mass makes the collective line invariant, but
without delay-layer left balance a transverse history can feed the
collective equation.  Thus the complete-history operator has the exact block
form
\[
 \cL_N=
 \begin{pmatrix}\cL_\parallel&C_N\\0&\cL_{\perp,N}\end{pmatrix},
\]
not in general a direct sum.  Since
\(\cL_{\perp,N}^{-1}=G_{\perp,N}\), the bounded codomain row operation
\begin{equation}
 T_N(y_\parallel,y_\perp)
 =\bigl(y_\parallel-C_NG_{\perp,N}y_\perp,y_\perp\bigr)
 \quad\hbox{satisfies}\quad
 T_N\cL_N=\cL_\parallel\oplus\cL_{\perp,N}.            \label{eq:triangular-row-reduction}
\end{equation}
Consequently the Fredholm index and kernel/cokernel dimensions are inherited
from the scalar block.  If \(\psi\) spans the scalar cokernel, normalized by
the collective complement \(e_N=(e_\parallel,0)\), then the full functional
is
\begin{equation}
 \Psi_N(y_\parallel,y_\perp)
 =\psi\bigl(y_\parallel-C_NG_{\perp,N}y_\perp\bigr).   \label{eq:triangular-cokernel}
\end{equation}
For a collective parameter or inhomogeneity
\((g_\parallel,0)\), \(\Psi_N(g_\parallel,0)=\psi(g_\parallel)\).
The normalized root, slope, and orientation are therefore exactly the scalar
ones whenever the scalar phase-fixed complete-history root and collective
preparation are supplied.

The triangular correction is uniform.  With
\[
 \delta_j=\frac12\norm{\pi^TB_j-\tfrac12\pi^T}_1,
 \qquad \delta_B=\delta_0+\delta_1,
\]
oscillation duality and the complete-line Green bound give
\begin{equation}
 \norm{C_NG_{\perp,N}}
 \le\frac{391}{2000}\delta_B\le\frac{391}{2000}.       \label{eq:triangular-row-bound}
\end{equation}
No delayed-history residence factor appears here: on complete-line sup
norms, a fixed time translation has norm one.  This is a statement about
complete-history matching.  Replacing the history jump by a current-state
projection would discard the delayed transverse correction in
\cref{eq:triangular-cokernel}.

For structural perturbations, exact splitting is no longer required.  We
briefly prove the scalar part of \cref{thm:one-gap}, since it explains the
specific radius in \cref{eq:rho-star}.  Put
\[
 d_N(s,\cR)=\widetilde d_N(\nu_{0,N}+s,\cR).
\]
For \(r=\norm{\cR}\), Taylor's formula and
\cref{eq:root-bounds} give, on
\(|s|\le2B_*r/m_*\),
\begin{equation}
 \abs{\partial_sd_N(s,\cR)-a_N}
 \le M_*\left(\abs s+r\right)
 \le M_*\left(1+\frac{2B_*}{m_*}\right)r
 \le\frac{m_*}{2}.                                    \label{eq:root-slope-strip}
\end{equation}
Thus the gap is strictly monotone with the orientation of \(a_N\).
Moreover, at \(s=\pm2B_*r/m_*\), the linear term dominates the residual,
so the endpoint signs are opposite.  The first two bounds in
\cref{eq:rho-star} keep these endpoints inside the common cylinder; the
third proves \cref{eq:root-slope-strip}; and the fourth keeps the
graph--trace range solve inside its validated ball.  Existence and
uniqueness follow.

Writing the root as \(s_N\) and expanding at \((0,0)\) yields
\[
 0=a_Ns_N+\ell_N[\cR]+\cE_N,\qquad
 \abs{\cE_N}\le\frac{M_*}{2}(\abs{s_N}+r)^2.
\]
The root enclosure then gives \cref{eq:root-remainder}.  This calculation
also shows why a small residual norm without a uniform simple-root bound is
insufficient: the displacement is divided by \(a_N\).

\subsection{The shared-resource history graph and selected traces}
\label{subsec:shared-proof}

We now prove the analytic steps used in \cref{thm:shared-resource}.  This
result has a different role from the preceding triangular transfer: it
constructs its own preparation-indexed planar histories for a heterogeneous
shared-resource network.  In particular, the graph and trace range problems
below are not additional unstated hypotheses of the theorem.
Here the ``invariant-tail levels'' in the admissible preparation are the two
zero levels \(\mathscr H_\alpha=0\) used below.  Arbitrary independently
chosen outer levels or outer RFDE slow manifolds are not included in this
preparation class.

\subsubsection{Exact chart and dimension-free model fit}

For a complete chart history \(\phi\), write
\[
 \cL_{N,\zeta}[\phi](s)
 =\sum_{k=0}^m(B_{k,N}+\zeta R_{k,N})
       \{\phi(s)-\phi(s-\theta_k)\}.
\]
Set
\[
 g_N=A_N^{-1}P_{\perp,N}c_N,
 \qquad Z_N(X,h)=g_NX^2+h.
\]
Substitution of the fold chart preceding \cref{eq:shared-gap} into
\cref{eq:shared-rfde-v,eq:shared-rfde-w} gives the following identities,
with no Taylor remainder:
\begin{align}
 X'={}&Y-\alpha X^2\notag\\
 &+\delta\left[-2X\pi_N^T\diag(c_N)Z_N
          -\frac\beta3X^3
          +K\pi_N^T\cL_{N,\zeta}[\mathbf1X]\right]\notag\\
 &+\delta^2\left[-\pi_N^T\diag(c_N)Z_N^{\circ2}
          +K\pi_N^T\cL_{N,\zeta}[Z_N]\right]
 -\delta^3\beta X\pi_N^TZ_N^{\circ2}
 -\frac{\delta^4\beta}{3}\pi_N^TZ_N^{\circ3},
                                                               \label{eq:shared-exact-X}\\
 Y'={}&-X+\delta\nu,                                  \label{eq:shared-exact-Y}\\
 \delta h'={}&A_Nh\notag\\
 &+\delta\left\{P_{\perp,N}\left[-2X\diag(c_N)Z_N
               +K\cL_{N,\zeta}[\mathbf1X]\right]
               -2g_NX X'\right\}\notag\\
 &+\delta^2P_{\perp,N}\left[-\diag(c_N)Z_N^{\circ2}
             -\beta X^2Z_N+K\cL_{N,\zeta}[Z_N]\right]
 -\delta^3\beta X P_{\perp,N}Z_N^{\circ2}
 -\frac{\delta^4\beta}{3}P_{\perp,N}Z_N^{\circ3}.
                                                               \label{eq:shared-exact-h}
\end{align}
The term \(-2g_NXX'\) is the derivative of the translation \(g_NX^2\).
The choice of \(g_N\) is forced by
\(A_Ng_NX^2=P_{\perp,N}c_NX^2\); hence the singular stable graph in these
coordinates is \(h=0\), and the reduced singular field is
\[
 q_0(X,Y)=(Y-\alpha X^2,-X)^T.
\]

We record why the constants in this exact chart do not hide a dependence on
the number of nodes.  The Poisson expansion and Dobrushin contraction give
\begin{equation}
 \norm{e^{A_Nt}}_{E_N\to E_N}\le e^{-D\gamma t},
 \qquad \norm{A_N^{-1}}_{E_N\to E_N}\le(D\gamma)^{-1}.
                                                               \label{eq:poisson-bound}
\end{equation}
Moreover, in the norm \cref{eq:shared-norm},
\begin{equation}
 \norm{P_{c,N}},\norm{P_{\perp,N}}\le1,\qquad
 \norm{x}_\infty\le\norm{x}_N,\qquad
 \norm{x_1\circ\cdots\circ x_j}_N
 \le3\prod_{\ell=1}^j\norm{x_\ell}_N.               \label{eq:shared-products}
\end{equation}
Indeed, subtracting \((\pi_N^Tx)\mathbf1\) does not change oscillation,
and a coordinate differs from the \(\pi_N\)-average by at most
\(\osc x\).  The last inequality follows by first estimating each
coordinate and then using \(\norm{y}_N\le3\norm{y}_\infty\).
Consequently \cref{eq:shared-curvature,eq:shared-layer-bounds},
\cref{eq:poisson-bound}, and \cref{eq:shared-products} bound every
multilinear derivative in \cref{eq:shared-exact-X,eq:shared-exact-h} by one
constant independent of \(N\) on a common chart box.

For later use, put
\[
 \cX_N=\R^2\times E_N,
 \qquad
 \cY_N=\cX_N^{m+2}
\]
with maximum product norms.  Current evaluation and the \(m+1\) fixed
delay evaluations are represented by the fixed atomic measure
\begin{equation}
 \mathbf M_N=J_{-1}\delta_0+\sum_{k=0}^mJ_k\delta_{-\theta_k},
 \qquad \norm{\mathbf M_N}_{\rm TV}\le m+2,           \label{eq:shared-evaluation-TV}
\end{equation}
where every insertion \(J_k\) has norm one.  The matrices depending on
\(\zeta\) remain in the polynomial nonlinearities, so the evaluation
measure is parameter independent.  This avoids differentiating moving
Dirac masses and makes its affine-history anchor derivatives identically
zero.

\begin{lemma}[Uniform graph and trace range]
\label{lem:shared-graph-trace}
Fix the common bounds in \cref{thm:shared-resource} and one bounded class of
admissible preparations.  After choosing the fixed exponent \(p>4\)
sufficiently large, there are \(\delta_*,\zeta_*,C,c>0\) and an integer
\(M\), all independent of \(N\) and of the preparation in that class, with
the following properties.

For \(0<\delta<\delta_*\), \(\nu\in I_\nu\), and
\(|\zeta|<\zeta_*\), let
\(S_\delta=\sqrt{2p\log(1/\delta)}\).  On a buffered tube about
\(\gamma_0([-S_\delta,S_\delta])\), the exact system
\cref{eq:shared-exact-X,eq:shared-exact-h} has an injective invariant
complete-history graph.  Its reduced field has the expansion
\begin{equation}
 Q_{N,\delta,\nu,\zeta}
 =q_0+\delta q_{1,N}+\delta^2q_{2,N}
       +\delta^3\mathcal R_{3,N},                     \label{eq:shared-graph-expansion}
\end{equation}
and, for the finite derivative list \(a\le3\), \(i\le1\), \(j\le2\),
\begin{equation}
 \abs{D_u^a\partial_\nu^i\partial_\zeta^j
       \mathcal R_{3,N}(\gamma_0(s)+\xi)}
 \le C\langle s\rangle^M e^{c|s|}                   \label{eq:shared-graph-remainder}
\end{equation}
on the nested retained tube.  The graph equation and its displayed mixed
derivatives have a right inverse bounded uniformly in \(N\).

The prepared tail conditions
\(\mathscr H_\alpha(z^a(-R))=\mathscr H_\alpha(z^r(R))=0\), where
\(R=S_\delta+B\) and \(B>2\Theta_*+2\) is fixed, together with the phase
\(X^a(0)=X^r(0)=0\), select unique one-sided traces.  Their retained central
pieces are exact histories on the unprepared graph, and their trace--phase--
endpoint linearization is onto with a uniformly bounded right inverse.
In the tame weighted norms used in the proof,
\begin{align}
 \norm{z^\sigma-\gamma_0}&\le C\delta,&
 \norm{\partial_\nu z^\sigma}&\le C\delta,\notag\\
 \norm{\partial_\zeta z^\sigma}
  +\norm{\partial_{\zeta\zeta}z^\sigma}&\le C\delta^2,
 &\sigma&\in\{a,r\}.                                  \label{eq:shared-trace-jets}
\end{align}
The same \(O(\delta^2)\) estimate holds for the mixed derivatives obtained
by applying one additional \(\nu\)-derivative to a \(\zeta\)-derivative.
\end{lemma}

\begin{proof}
We include the construction because it is the range bridge between the
finite-dimensional coefficient calculation and the complete-history root.
In the coordinates
\[
 \chi=-2\alpha X,
 \qquad d=Y-\alpha X^2+\frac1{2\alpha},
\]
prepare \(q_0\) by one common compact cutoff which is the identity on the
continuous depth-two \(q_0\)-flow hull of the retained tube.  Prepare every
current and delayed stable slot by the same componentwise smooth saturation
followed by \(P_{\perp,N}\).  Estimate \cref{eq:shared-products} shows that
all derivatives of these preparations are bounded, without a factor of
\(N\), by a polynomial \(P(S_\delta)\).  The preparation is the physical
system on a common stable ball containing the entire retained depth-two
history hull.

For a bounded planar field \(Q\) and a graph \(H:\R^2\to E_N\), let
\[
 \mathcal I_{Q,H}(u)(\vartheta)
 =\bigl(\Phi_Q^\vartheta u,H(\Phi_Q^\vartheta u)\bigr),
 \qquad -\Theta_*\le\vartheta\le0.
\]
For each target \(\delta\), freeze this preparation at
\(S=S_\delta\), introduce a dummy amplitude
\(\rho\in[-\delta,\delta]\), and recover the physical chart at
\(\rho=\delta\).  Writing the prepared version of
\cref{eq:shared-exact-X,eq:shared-exact-h} as
\(u'=q_0(u)+\rho F_N\),
\(\rho h'=A_Nh+\rho G_N\), its special-flow transform is
\begin{align}
 \mathcal T_{Q,N}(Q,H)(u)
 &=q_0(u)+\rho F_N
   \bigl(u,H(u),\mathbf M_N\mathcal I_{Q,H}(u);\rho,\nu,\zeta\bigr),
                                                               \label{eq:shared-TQ}\\
 \mathcal T_{H,N}(Q,H)(u)
 &=\rho\int_0^\infty e^{A_Nr}G_N
   \bigl(\mathcal E_{Q,H}(\Phi_Q^{-\rho r}u);
               \rho,\nu,\zeta\bigr)\,\dd r.          \label{eq:shared-TH}
\end{align}
Here \(\mathcal E_{Q,H}\) denotes the current state and the evaluation
tuple in \cref{eq:shared-evaluation-TV}.  A reduced-flow variation over a
time \(\Theta_*+\delta r\), followed by
\cref{eq:poisson-bound,eq:shared-evaluation-TV}, gives the common
majorant
\begin{equation}
 C\delta P(S_\delta)e^{cS_\delta}
 \int_0^\infty(1+r^M)
 e^{-\{D\gamma-c\delta(1+S_\delta)\}r}\,\dd r=o(1)   \label{eq:shared-graph-contraction}
\end{equation}
for the Lipschitz constant of \((\mathcal T_{Q,N},\mathcal T_{H,N})\).
The same estimate holds for the highest current block after any derivative
in the finite list of the lemma; all other terms contain already controlled
lower blocks.  More explicitly, order
\[
 J_{a,b,i,j}=D_u^a\partial_\rho^b\partial_\nu^i
                 \partial_\zeta^j,\qquad
 b\le3,\quad i\le1,\quad j\le2,
\]
first by total order and then by parameter order.  Fa\`a di Bruno's formula
and the flow variational equations put each current block in the form
\begin{equation}
 \mathbf J\mathcal T_N(Z)
 =B_{\mathbf J}(\mathbf J_{<}Z)
 +\rho\,\mathcal A_{\mathbf J}(\mathbf J_{<}Z)
          [\mathbf JZ].                               \label{eq:shared-jet-triangle}
\end{equation}
If a derivative hits a flow argument, the vector-field jet gains one state
slot and loses one parameter derivative, so it precedes the current block;
a product of two positive-order jets has two lower total orders.  Derivatives
of \(\mathbf M_N\) vanish, and
\cref{eq:shared-evaluation-TV} bounds every remaining history evaluation.
The coefficient of the last term in
\cref{eq:shared-jet-triangle} has the same \(o(1)\) bound as
\cref{eq:shared-graph-contraction}.  Induction over this finite block list
gives invariant jet balls and, for consecutive Picard iterates, recurrences
of the form
\[
 a_{n+1}\le\kappa a_n+Cq^n,\qquad \max\{\kappa,q\}<1,
\]
with the same constants for every \(N\).  Uniform convergence of a jet and
its first derivative, followed by the fundamental theorem of calculus on
state and parameter line segments, identifies the limits as genuine mixed
derivatives.  This proves convergence of the mixed Picard jets, uniformly
for \(|\rho|\le\delta\), and gives
\begin{equation}
 \norm{(I-D\mathcal T_N)^{-1}}\le2                   \label{eq:shared-graph-range}
\end{equation}
after decreasing \(\delta_*\).  Taylor's formula in \(\rho\), evaluated at
\(\rho=\delta\), gives
\cref{eq:shared-graph-expansion}; the corresponding graph-pair remainder is
the Banach-valued integral
\[
 \widehat{\mathcal R}_{3,N}(\rho)
 =\frac12\int_0^1(1-t)^2
       \partial_\rho^3(Q_{N,t\rho},H_{N,t\rho})\,\dd t.
\]
The \(Q\)-component of \(\widehat{\mathcal R}_{3,N}\) is
\(\mathcal R_{3,N}\) in \cref{eq:shared-graph-expansion}.

For the pointwise estimate, split the stable convolution at
\(L\log(1/\delta)\).  Decrease \(\delta_*\) so that
\(c\delta(1+S_\delta)<D\gamma/2\), and then choose \(L\), independently of
\(N\), so that
\begin{equation}
 CP(S_\delta)e^{cS_\delta}
 \{1+L\log(1/\delta)\}^M
 e^{-D\gamma L\log(1/\delta)/2}=O(\delta^6).          \label{eq:shared-convolution-tail}
\end{equation}
On the finite part, the additional reduced-flow backtrack is at most
\(\delta L\log(1/\delta)=o(1)\).  The same triangular induction in nested
normal tubes now gives \cref{eq:shared-graph-remainder}.  At
\(\rho=0\), one derivative uses one delay-support backtrack; two derivatives
use at most two backtracks and the finite semigroup moments
\(\int_0^\infty r^ke^{A_Nr}\dd r\), \(k=0,1\), which are bounded by
\(k!(D\gamma)^{-k-1}\).  Thus the first two coefficients use only the
declared depth-two uncut hull.  Finally,
the fixed-point identity in \cref{eq:shared-TH}, split at any time, is the
mild stable equation.  Present-time evaluation is a left inverse, so the
resulting complete-history map is injective and its retained histories solve
the unprepared RFDE exactly.

For the trace solve, repeat the graph construction with the frozen target
\(S_\delta+B\).  Extend its reduced field in the normal coordinate by one
fixed finite-order extension: it agrees with the graph field on the
shrinking retained tube and equals \(q_0\) near the two endpoints
\(\gamma_0(\pm R)\).  The extension uses the unscaled normal distance from
the boundary of the shrinking tube, so it introduces no inverse power of
\(\delta\) into the state or parameter estimates.  Only the central pieces
which bootstrap back into the exact graph are called RFDE histories; the
outer planar tails are selection devices.

It remains to solve for the two selected planar traces.  Along
\[
 \gamma_0(s)=\left(-\frac{s}{2\alpha},
                    \frac{s^2-2}{4\alpha}\right),
\]
the linearized operator is
\[
 L_0(U,V)=(U'-sU-V,V'+U).
\]
Put \(E(s)=e^{s^2/2}\),
\(I(s)=\int_0^se^{t^2/2}\,\dd t\), and
\[
 \tau(s)=(-1,s)^T,
 \qquad n(s)=(I(s),E(s)-sI(s))^T.
\]
These are respectively the tangent and Gaussian normal solutions, with
\(\det[\tau,n]=-E\).  If \(\xi=a\tau+bn\) and \(L_0\xi=f\), then
\begin{align}
 a'&=E^{-1}\{-(E-sI)f_1+If_2\}=:A_f,&
 b'&=E^{-1}(sf_1+f_2)=:B_f.                           \label{eq:shared-ab}
\end{align}
The phase \(U(0)=0\) is exactly \(a(0)=0\).  With the no-incoming normal
coordinate prescribed at the outer endpoint, the two one-sided inverses are
\begin{align}
 a_a(s)&=-\int_s^0A_f(t)\,\dd t,&
 b_a(s)&=\beta_a+\int_{-R}^sB_f(t)\,\dd t,\notag\\
 a_r(s)&= \int_0^sA_f(t)\,\dd t,&
 b_r(s)&=\beta_r-\int_s^RB_f(t)\,\dd t.               \label{eq:shared-one-sided-green}
\end{align}
The elementary estimates
\[
 |I(s)|\le CE(s)\langle s\rangle^{-1},
 \qquad |E(s)-sI(s)|\le CE(s)
\]
and the one-sided Gaussian tail bounds show that, for fixed \(c_0,M_0\),
\begin{equation}
 \norm{\mathcal G_R^\sigma(f,\beta_\sigma)}_{c_0,M_0+d}
 \le C\left(\norm f_{c_0,M_0}
 +e^{R^2/2-c_0R}\langle R\rangle^{-M_0}|\beta_\sigma|\right),
                                                               \label{eq:shared-green-bound}
\end{equation}
where \(d\le4\), \(C\) is independent of \(R\), and
\(\norm g_{c,M}=\sup e^{-c|s|}\langle s\rangle^{-M}|g(s)|\).
Thus the Gaussian growth is paired with its one-sided integral before a
norm is taken.

The tail equation \(\mathscr H_\alpha=0\) is not an assumed stable
manifold.  At an endpoint it is exactly
\begin{equation}
 E(s)b=\alpha\{-a+I(s)b\}^2,                           \label{eq:shared-tail-level}
\end{equation}
so the implicit-function theorem solves \(b=\mathcal B_\sigma(a)\) with
zero value and derivative at the origin, uniformly in the normalized
boundary norm of \cref{eq:shared-green-bound}.  For
\(z=\gamma_0+\xi\), the nonlinear trace equation has the form
\(L_0\xi=N_{\delta,\nu,\zeta}(s,\xi)\), with
\[
 |N(s,0)|\le C\delta e^{c|s|}\langle s\rangle^M,
 \qquad
 |D_\xi N(s,\xi)|\le
 C\{\delta e^{c|s|}\langle s\rangle^M+|\xi|\}.
\]
Use two tame weights separated by the finite loss \(d\) in
\cref{eq:shared-green-bound}.  On a ball of radius \(C\delta\), the
one-sided fixed-point map defined by
\cref{eq:shared-one-sided-green,eq:shared-tail-level} has Lipschitz factor
\begin{equation}
 C\delta e^{cR}\langle R\rangle^M=o(1),
 \qquad R=S_\delta+B.                                 \label{eq:shared-trace-contraction}
\end{equation}
It therefore has a unique fixed point, and its linearization has inverse
norm at most two.  The \(\nu\)-source is \(O(\delta)\); atomwise row
neutrality makes the first \(\zeta\)-source \(O(\delta^2)\).
Differentiating the same fixed-point equation proves
\cref{eq:shared-trace-jets} and the stated mixed bounds.

The graph is solved first and the trace equations second.  Their combined
derivative is block lower triangular; \cref{eq:shared-graph-range} and the
last trace inverse therefore give a uniform right inverse for the full
graph--trace--phase--endpoint range problem.  Finally, in the normal
coordinate one has the exact identity
\[
 d(\gamma_0(s)+(U,V))=V+sU-\alpha U^2.
\]
The first estimate in \cref{eq:shared-trace-jets} makes this
\(O(\delta e^{cS_\delta}\langle S_\delta\rangle^{M+1})\), which is
\(o(\delta^\kappa)\) for every fixed \(\kappa<1\).  Hence the central trace
and all of its retained delay backtracks lie inside the uncut shrinking
graph, completing the exact-history assertion.
\end{proof}

The range assertion in \cref{lem:shared-graph-trace} is the one required by
this selected object: the invariant-graph residual followed by its tangent
one-sided trace problem.  It is not an inverse for arbitrary forcing off the
graph in the ambient RFDE history space, and no such stronger inverse is
used below.

\subsubsection{The mixed curvature return and the normalized gap}

\begin{proof}[Proof of \cref{thm:shared-resource}]
We calculate the coefficient after, rather than before, the range problem
has been solved.  Let \(H_{1,N}\) be the first stable graph coefficient in
the expansion furnished by \cref{lem:shared-graph-trace}.  On the singular
orbit,
\begin{equation}
 X_0(s)-X_0(s-\theta_k)=-\frac{\theta_k}{2\alpha}.
                                                               \label{eq:shared-affine-delay}
\end{equation}
Differentiating the order-\(\delta\) stable graph equation and using
\cref{eq:shared-affine-delay} gives
\[
 A_ND_\zeta H_{1,N}(\gamma_0(s))
 -\frac{K}{2\alpha}P_{\perp,N}\dot M_{1,N}\mathbf1=0.
\]
Thus
\begin{equation}
 D_\zeta H_{1,N}(\gamma_0(s))=h_{*,N},\qquad
 h_{*,N}=\frac{K}{2\alpha}A_N^{-1}P_{\perp,N}
             \dot M_{1,N}\mathbf1.                  \label{eq:hstar-rn-a}
\end{equation}
This vector is independent of \(s\).  Put
\begin{equation}
 r_N=\pi_N^T\diag(c_N)h_{*,N}.                         \label{eq:hstar-rn}
\end{equation}
Because \(\pi_N^TR_{k,N}=0\) as a row identity, the structural derivative
of \(\pi_N^T\cL_{N,\zeta}[\phi]\) vanishes for every vector history
\(\phi\).  It therefore removes the direct critical delay term, the
translated history \(g_NX^2\), and the flow-history corrections before any
moment reduction is made.  The translation is independent of \(\zeta\).
The only order-\(\delta^2\) structural term left in
\cref{eq:shared-exact-X} is
\(-2X\pi_N^T\diag(c_N)h_{*,N}\).  Consequently
\begin{equation}
 D_\zeta q_{1,N}=0,\qquad
 D_\zeta q_{2,N}(\gamma_0(s),\nu,0)
   =\left(\frac{r_N}{\alpha}s,0\right)^T.              \label{eq:q2-return}
\end{equation}
This proves both the cancellation of the lower-order response and the
mixed-curvature return; it is an identity for the actual graph jets of
\cref{eq:shared-graph-expansion}, not a formal current-state expansion.

We next include the trace and endpoint contributions.  The normalized
adjoint and the first-integral differential along \(\gamma_0\) are
\begin{equation}
 \psi(s)=e^{-s^2/2}(s,1)^T,\qquad
 \nabla\mathscr H_\alpha(\gamma_0(s))
       =\frac{\alpha e}{2}\psi(s).                    \label{eq:shared-adjoint}
\end{equation}
The first reduced coefficient is
\begin{equation}
 q_{1,N,X}=-\kappa_NX^3
          +K\pi_N^T\cL_{N,0}[\mathbf1X],\qquad
 q_{1,N,Y}=\nu.                                       \label{eq:shared-q1}
\end{equation}
Along \(\gamma_0\), the delay term in \cref{eq:shared-q1} is constant in
\(s\), by \cref{eq:shared-affine-delay}, and hence pairs to zero with the
odd function \(se^{-s^2/2}\).  Using
\(\int e^{-s^2/2}\dd s=\sqrt{2\pi}\),
\(\int s^2e^{-s^2/2}\dd s=\sqrt{2\pi}\), and
\(\int s^4e^{-s^2/2}\dd s=3\sqrt{2\pi}\), we obtain
\begin{align}
 \int_\R\psi^Tq_{1,N}(\gamma_0(s),\nu)\,\dd s
 &=\sqrt{2\pi}\left(\nu+\frac{3\kappa_N}{8\alpha^3}\right)
 =\sqrt{2\pi}(\nu-\nu_{0,N}),                        \label{eq:shared-leading-gap}\\
 \int_\R\psi^TD_\zeta q_{2,N}(\gamma_0(s),\nu,0)\,\dd s
 &=\sqrt{2\pi}\frac{r_N}{\alpha}.                    \label{eq:shared-melnikov}
\end{align}

For clarity, these whole-line integrals do not replace a finite-section
boundary-value problem.  The two endpoint values of
\(\mathscr H_\alpha\) are zero by construction, and integrating
\(\dd\mathscr H_\alpha/\dd s\) along the two one-sided traces includes the
phase, both endpoints, and the moving complete histories.  The Green bounds
in \cref{lem:shared-graph-trace} and
\cref{eq:shared-graph-remainder} allow the traces to be replaced by
\(\gamma_0\) in the leading terms.  The omitted Gaussian tails and the
difference between two preparations in the declared class are bounded,
with the same parameter derivatives, by
\begin{equation}
 C e^{-S_\delta^2/2+cS_\delta}
       \langle S_\delta\rangle^M
 =C\delta^{p-o(1)}.                                   \label{eq:shared-gaussian-tail}
\end{equation}
Choosing the fixed \(p\) after the finite tame losses makes this smaller
than the algebraic remainders below.  The pairings in
\cref{eq:shared-leading-gap,eq:shared-melnikov} therefore yield, uniformly
in \(N\), in the preparation class, and on the common parameter box,
\begin{align}
 \frac{\mathscr D_{N,\mathcal P_N}}\delta
 &=\sqrt{2\pi}(\nu-\nu_{0,N})+O(\delta),
                                                               \label{eq:shared-gap-expansion}\\
 \partial_\nu\mathscr D_{N,\mathcal P_N}
 &=\delta\sqrt{2\pi}+O(\delta^2),                    \label{eq:shared-gap-nu}\\
 \partial_\zeta\mathscr D_{N,\mathcal P_N}
 &=\delta^2\sqrt{2\pi}\frac{r_N}{\alpha}
       +O(\delta^3+\delta^2|\zeta|),                 \label{eq:shared-gap-zeta}\\
 \partial_{\zeta\zeta}\mathscr D_{N,\mathcal P_N}
 &=O(\delta^2).                                       \label{eq:shared-gap-zz}
\end{align}
Thus the graph, trace, endpoint, and remainder estimates needed for the
root are all part of the same uniform construction.

Choose a common \(c_0>0\) smaller than the boundary distance of the
\(\nu_{0,N}\) from \(\partial I_\nu\).
\Cref{eq:shared-gap-expansion,eq:shared-gap-nu} give opposite endpoint
signs and strict monotonicity on
\(|\nu-\nu_{0,N}|<c_0\) for small \(\delta,\zeta\), proving the unique root
in \cref{eq:shared-root}.  Along that root,
\begin{equation}
 \partial_\zeta\nu_{c,N,\mathcal P_N}
 =-\delta\frac{r_N}{\alpha}
       +O(\delta^2+\delta|\zeta|).                    \label{eq:shared-root-derivative}
\end{equation}
Integrating \cref{eq:shared-root-derivative} from \(0\) to \(\zeta\) and
multiplying by \(\delta^2\) gives
\cref{eq:shared-response,eq:shared-coeff}, with a constant independent of
\(N\) and of the declared preparation.  On the section \(X=0\), the map
\(Y\mapsto\mathscr H_\alpha(0,Y)\) has derivative \(\alpha e/2\ne0\) at
the singular crossing.  Hence a zero of the gap is equality of the two
reduced current states.  Injectivity of the history graph in
\cref{lem:shared-graph-trace} then makes it equality of the two retained
complete RFDE histories.  The common value of \(\alpha\) and
\cref{eq:shared-nondegeneracy} also give
\(\inf_N|\mathscr C_N|=\alpha^{-1}\inf_N|r_N|>0\), as asserted.
This completes the proof of \cref{thm:shared-resource}.
\end{proof}

\subsubsection{A uniform nonzero family without a synchrony quotient}

It remains to verify that the nondegeneracy hypothesis is nonempty without
replicating a lower-dimensional synchrony class.  Use the explicit profile
stated after \cref{thm:shared-resource}.  Direct summation gives
\begin{equation}
 \pi_N^Ts_N=0,\qquad \pi_N^Ts_N^{\circ2}=1,\qquad
 \norm{s_N}_\infty
 =\sqrt{\frac{3(N-1)}{N+1}}<\sqrt3,                  \label{eq:shared-profile-audit}
\end{equation}
and its entries are pairwise distinct.  Thus
\(c_N=\mathbf1+\sigma s_N\) is positive for
\(0<\sigma<1/\sqrt3\).  Moreover,
\[
 \norm{s_N\pi_N^Tx}_N
 \le2\sqrt3\,\norm{x}_N,
\]
so the layer bound is uniform.  Since
\(B_{k,N}\pm\zeta s_N\pi_N^T\) has entries
\((b\pm\zeta s_{i,N})/N\), both perturbed layers are positive throughout
the common interval \(|\zeta|<b/\sqrt3\).

Here \(A_N=-D\Id\) on \(E_N\) and
\(\dot M_{1,N}\mathbf1=(\theta_0-\theta_1)s_N\).
Together, \cref{eq:hstar-rn-a,eq:hstar-rn,eq:shared-profile-audit} give
\begin{equation}
 r_N=-\frac{K\sigma(\theta_0-\theta_1)}{2D},
 \qquad
 \mathscr C_N=\frac{K\sigma(\theta_0-\theta_1)}{2D}\ne0,
                                                               \label{eq:constant-witness}
\end{equation}
the same number for every \(N\ge2\).

Finally, suppose a synchrony polydiagonal had a cell containing distinct
nodes \(i\ne j\).  Test invariance on a time-constant history with
\(v_i=v_j\ne1\).  The balanced delayed term vanishes, the two rows of
\(P_{c,N}\) give identical instantaneous coupling, and the shared resource
and cubic terms agree at the two nodes.  Their quadratic terms differ by
\((c_{i,N}-c_{j,N})(v_i-1)^2\ne0\).  The vector field is therefore not
tangent to that polydiagonal.  Since all curvatures are distinct, only the
singleton partition survives.  Thus \cref{eq:constant-witness} is a genuine
transverse heterogeneous-network response and not a lifted two-node root.

\subsection{Nonlinear synchronization and collective defect}

The same diameter estimate applies to the nonlinear network while all node
voltages remain in the strip.  With \(M_0\) the retained-history diameter at
entrance, the validated nonlinear comparison gives
\begin{equation}
 M(t)\le M_0e^{-(t-t_0)/10}.                            \label{eq:nonlinear-sync}
\end{equation}
Let \((\bar v,\bar w)=(\pi^Tv,\pi^Tw)\) and write
\(\pi^TB_j=\frac12\pi^T+b_j^T\).  Row mass makes
\(b_j^T\mathbf1=0\), so oscillation duality controls the delayed
left-imbalance.  The collective scalar equation has a forcing satisfying
\begin{align}
 |R_{\rm coll}(t)|\le{}&
 \frac{391}{20000}
 \{\delta_0M(t-\tau_0)+\delta_1M(t-\tau_1)\}\notag\\
 &+\frac{1403}{400}M(t)^2
 +\frac3{800}\{M(t-\tau_0)^2+M(t-\tau_1)^2\}.          \label{eq:collective-defect-resolved}
 \end{align}
With \(\mathcal H_M(t)=\sup_{t-r\le s\le t}M(s)\), this gives
 \begin{equation}
 |R_{\rm coll}(t)|
 \le\frac{391}{20000}\delta_B\mathcal H_M(t)
    +\frac{703}{200}\mathcal H_M(t)^2.                 \label{eq:collective-defect-envelope}
 \end{equation}
Retaining the prescribed history for the full delayed residence intervals
gives the resolved forward bound
 \begin{align}
 \norm{R_{\rm coll}}_{L^1(t_0,\infty)}
 \le{}&\frac{391}{20000}
 \{\delta_0(10+4\sqrt5)+\delta_1(10+5\sqrt5)\}M_0\notag\\
 &+\left(\frac{703}{40}+\frac{27\sqrt5}{800}\right)M_0^2,\label{eq:collective-defect-resolved-L1}\\
 \intertext{and, when only \(\delta_B\) is retained,}
 \norm{R_{\rm coll}}_{L^1(t_0,\infty)}
 \le{}&\frac{391(2+\sqrt5)}{4000}\delta_BM_0
       +\frac{56483}{3200}M_0^2.                       \label{eq:collective-defect}
 \end{align}
These constants are independent of \(N\) and of the admitted topology.
The term \(27\sqrt5/800\) is the exact residence correction for the two
delayed quadratic remainders, while the factors \(10+\tau_j\) retain the
corresponding delayed linear histories.  They cannot be removed for an
arbitrary retained history.

For the left-balanced subclass \(\delta_B=0\), every linear collective term
cancels and \cref{eq:collective-defect-resolved} reduces to
\begin{align}
 |R_{\rm coll}(t)|\le{}&\frac{1403}{400}M(t)^2
 +\frac3{800}\{M(t-\tau_0)^2+M(t-\tau_1)^2\},\notag\\
 \intertext{and hence, with
 \(\mathcal H_M(t)=\sup_{t-r\le s\le t}M(s)\),}
 |R_{\rm coll}(t)|\le{}&\frac{703}{200}\mathcal H_M(t)^2,\notag\\
 \norm{R_{\rm coll}}_{L^1(t_0,\infty)}
 \le{}&\left(\frac{703}{40}+\frac{27\sqrt5}{800}\right)M_0^2
 \le\frac{56483}{3200}M_0^2.                           \label{eq:balanced-collective-defect}
\end{align}
Thus the proved uniform estimate for the enlarged class contains a resolved
term linear in \(M_0\).  Without an additional cancellation, the present
argument does not yield a purely quadratic enlarged-class estimate.

None of these estimates is, by itself, a basin theorem.  In the balanced
subclass, if a scalar route accepts entrance error
\(\delta_{\rm mean}\) and additive forcing whenever
\[
 L_0\delta_{\rm mean}+\norm{e}_{L^1}<\eta_{\rm route},
\]
and if \(d_{\rm strip}\) and \(d_{\rm lift}\) are respectively a scalar
strip margin and a product-basin lift margin, then the exact sufficient
network budget is
\begin{equation}
 R_0<\min\{d_{\rm strip},d_{\rm lift}\},\qquad
 L_0R_0+\frac{56483}{3200}R_0^2<\eta_{\rm route}.       \label{eq:async-budget}
\end{equation}
Here \(R_0=\max\{\delta_{\rm mean},M_0\}\).  A first-strip-exit
contradiction proves that the hypotheses bootstrap until routing is
complete.  The scalar route, strip, lift, and gap constants in
\cref{eq:async-budget} remain model-specific release data.  No concrete
radius has been validated, and the released nonbalanced theorem does not
promote an asynchronous threshold or basin.

\subsection{Scope of the network theorem}

The leaky transverse and triangular-transfer theorem allows arbitrary finite
size, a nonnegative row-stochastic current layer with
\(\tau(Q)\le1/2\), and fixed nonnegative delayed layers of row mass \(1/2\).
It does not require common delay-layer left balance.  It remains conditional
on the scalar phase-fixed root and collective preparation, and it does not
include signed coupling, moving delay support, closing Dobrushin gaps,
heterogeneous node vector fields, an invariant asynchronous strip, or
physical onset.

The separate shared-resource structural-root theorem allows heterogeneous
curvature under a common Dobrushin gap, fixed delay support, an atomwise
row-neutral structural direction, and the bounded admissible preparation
class stated in \cref{thm:shared-resource}.  The graph and selected-trace
range solves for that class are constructed in
\cref{lem:shared-graph-trace}; they are not extra hypotheses.  The exact
finite-\(\delta\) root remains preparation indexed.  The theorem does not
cover order-one changes of the base generator, moving delay atoms, a closing
mixing gap, a preparation-independent root, or identification with an
independently selected physical outer canard.  These are mathematical
boundaries, not numerical omissions.
